% ============================================================
%  Chapter 11 --- Web Interface
% ============================================================
\chapter{Web Interface}
\label{ch:web-interface}

\section{Architecture Overview}

The \shiftdb{} web interface follows a classic \textbf{client-server} architecture:

\begin{figure}[H]
\centering
\begin{tikzpicture}[node distance=1.2cm and 2.5cm]
  \node[component, minimum width=4cm, fill=shiftPrimary!12] (browser) {React SPA\\(Vite + TypeScript)};
  \node[component, minimum width=4cm, right=of browser] (http) {HTTP/JSON};
  \node[component, minimum width=4cm, right=of http] (server) {C++ HTTP Server\\(raw sockets)};

  \draw[arr] (browser) -- (http) node[midway, above, font=\scriptsize] {REST calls};
  \draw[arr] (http) -- (server);
\end{tikzpicture}
\caption{Client-server communication via REST API.}
\end{figure}

\section{REST API Design}
\label{sec:rest-api}

\begin{table}[H]
\centering
\caption{Complete REST API reference.}
\label{tab:rest-api}
\begin{tabular}{l l l p{6cm}}
\toprule
\textbf{Method} & \textbf{Endpoint} & \textbf{Handler} & \textbf{Description} \\
\midrule
\code{POST} & \code{/api/query}            & \code{handleQuery}          & Execute SQL; body: \code{\{"sql":"..."\}} \\
\code{GET}  & \code{/api/schema}           & \code{handleGetSchema}      & All table schemas with columns, indexes, row counts \\
\code{GET}  & \code{/api/schema/:table}    & \code{handleGetTableSchema} & Single table schema \\
\code{GET}  & \code{/api/stats}            & \code{handleGetStats}       & Buffer pool + query metrics \\
\code{GET}  & \code{/api/history}          & \code{handleGetHistory}     & Last 100 queries with timing \\
\code{GET}  & \code{/api/tables}           & \code{handleGetTables}      & Simple list of table names \\
\code{GET}  & \code{/api/export/:table}    & \code{handleExport}         & Export as JSON or CSV \\
\code{POST} & \code{/api/import}           & \code{handleImport}         & Bulk import JSON data \\
\code{GET}  & \code{/api/health}           & (inline)                    & Health check: \code{\{"status":"ok"\}} \\
\bottomrule
\end{tabular}
\end{table}

\subsection{Query Execution Flow}

\begin{lstlisting}[style=cpp, caption={Simplified query handler.}]
HttpResponse handleQuery(const HttpRequest& req) {
    auto body = nlohmann::json::parse(req.body);
    std::string sql = body.value("sql", "");
    auto result = executor_.execute(sql);
    HttpResponse res;
    res.body = result.toJson().dump();
    return res;
}
\end{lstlisting}

The response JSON includes:
\begin{itemize}[nosep]
  \item \code{success} (bool), \code{message} (string)
  \item \code{columns} --- array of \code{\{name, type\}} objects
  \item \code{rows} --- 2D array of values
  \item \code{rowsAffected}, \code{executionTimeMs}
  \item \code{queryPlan} (for EXPLAIN)
\end{itemize}

\section{C++ HTTP Server Implementation}
\label{sec:http-server}

\shiftdb{}'s HTTP server is built from scratch using BSD sockets --- \textbf{zero external dependencies}.

\subsection{Server Lifecycle}

\begin{enumerate}[nosep]
  \item Create TCP socket (\code{socket()}).
  \item Bind to port (\code{bind()}) with \code{SO\_REUSEADDR}.
  \item Listen for connections (\code{listen()}).
  \item Accept loop: for each connection, spawn a \code{std::thread} to handle it.
  \item Parse HTTP request (method, path, headers, body).
  \item Route to handler, generate response.
  \item Send response and close socket.
\end{enumerate}

\subsection{Cross-Platform Support}

\begin{lstlisting}[style=cpp, caption={Platform-specific socket abstraction.}]
#ifdef _WIN32
  #include <winsock2.h>
  using socket_t = SOCKET;
  #define CLOSE_SOCKET closesocket
#else
  #include <sys/socket.h>
  using socket_t = int;
  #define CLOSE_SOCKET close
#endif
\end{lstlisting}

\subsection{CORS Handling}

The server adds CORS headers to every response, enabling the React frontend (running on a different port) to communicate:

\begin{lstlisting}[style=cpp, caption={CORS headers in response.}, numbers=none]
Access-Control-Allow-Origin: *
Access-Control-Allow-Methods: GET, POST, PUT, DELETE, OPTIONS
Access-Control-Allow-Headers: Content-Type, Authorization
\end{lstlisting}

\section{React Frontend Architecture}

\subsection{State Management}

The frontend uses \textbf{Zustand} for global state management --- a lightweight alternative to Redux:

\begin{table}[H]
\centering
\caption{Zustand store slices.}
\begin{tabular}{l p{12cm}}
\toprule
\textbf{State} & \textbf{Contents} \\
\midrule
\code{query}   & Current SQL text, execution results, loading state \\
\code{schema}  & Table schemas, column info, index metadata \\
\code{history} & Past queries with timing and success status \\
\code{stats}   & Buffer pool metrics, table statistics \\
\code{ui}      & Active view, sidebar state, theme \\
\bottomrule
\end{tabular}
\end{table}

\subsection{UI Components}

\begin{figure}[H]
\centering
\begin{tikzpicture}[
  node distance=0.5cm and 0.3cm,
  every node/.style={rectangle, draw=shiftPrimary!50, fill=shiftPrimary!5,
    rounded corners=2pt, font=\scriptsize, minimum height=0.6cm, align=center}
]
  \node[minimum width=11cm, minimum height=0.8cm] (app) {App.tsx (Layout Manager)};

  \node[minimum width=2.5cm, below left=0.6cm and 2cm of app] (sidebar) {Sidebar\\(Navigation)};
  \node[minimum width=7cm, below right=0.6cm and -2cm of app] (main) {Main Panel};

  \node[minimum width=3.2cm, below left=0.5cm and 0.8cm of main] (editor)  {SqlEditor\\(Monaco)};
  \node[minimum width=3.2cm, below right=0.5cm and 0.8cm of main] (results) {DataGrid\\(Results)};

  \node[minimum width=2.4cm, below=0.5cm of editor] (schema) {Schema\\Explorer};
  \node[minimum width=2.4cm, right=0.3cm of schema] (er) {ER\\Diagram};
  \node[minimum width=2.4cm, right=0.3cm of er] (qb) {Query\\Builder};
  \node[minimum width=2.4cm, right=0.3cm of qb] (mon) {Monitoring};

  \draw[arr] (app) -- (sidebar);
  \draw[arr] (app) -- (main);
  \draw[arr] (main) -- (editor);
  \draw[arr] (main) -- (results);
\end{tikzpicture}
\caption{Frontend component hierarchy.}
\label{fig:fe-hierarchy}
\end{figure}

\section{Node.js Backend (Alternative)}

For deployment to serverless platforms (Vercel), \shiftdb{} also provides a Node.js backend using \textbf{Express} + \textbf{better-sqlite3}:

\begin{itemize}[nosep]
  \item \code{server.cjs} implements the same REST API endpoints.
  \item Uses \code{sql.js} (SQLite compiled to WASM) for serverless compatibility.
  \item Schema, history, and stats endpoints mirror the C++ server's responses.
\end{itemize}

\begin{interviewbox}[Interview Corner]
\textbf{Q:} Why build a custom HTTP server instead of using a library?\\[4pt]
\textbf{A:} Zero external dependencies --- the C++ backend compiles without any third-party HTTP library.  This simplifies the build system and demonstrates low-level networking.  In production, you would use a battle-tested library (Boost.Beast, cpp-httplib) for robustness.

\vspace{6pt}
\textbf{Q:} What is CORS and why is it needed here?\\[4pt]
\textbf{A:} CORS (Cross-Origin Resource Sharing) is a browser security mechanism.  The React dev server runs on port 5173 but the API is on port 8080 --- different origins.  Without CORS headers, the browser blocks the API calls.

\vspace{6pt}
\textbf{Q:} How does the frontend maintain state?\\[4pt]
\textbf{A:} Using Zustand, a minimal state management library.  State slices (query, schema, history, stats, ui) are updated via actions that call the REST API and store responses.
\end{interviewbox}
